\subsection{半单李代数的嘉当子代数}

定理 2. 假定 k 的特征数为 0。设 g 是一个半单李代数，h 是 g 的嘉当子代数。那么 h 交换，且它的所有元素在 g 中半单（第一章，§6，第 3 节，定义 3）。

由于 \(\mathfrak{h} = \mathfrak{g}^0 (\mathfrak{h})\) ，\(\mathfrak{h}\) 约化（§1，命题 11），又因它幂零故交换。另一方面，\(\mathfrak{g}\) 的伴随表示在 \(\mathfrak{h}\) 上的限制半单（同处所引），故 \(\mathfrak{h}\) 的元素在 \(\mathfrak{g}\) 中半单（《代数》，第八章，§5，第 1 节）。

推论 1. 若 \(x \in \mathfrak{h}\) 且 \(y \in \mathfrak{g}^{\lambda}(\mathfrak{h})\) ，我们有 \([x, y] = \lambda(x)y\) 。

事实上，\(\mathfrak{g}^{\lambda (x)}(x) = \mathfrak{g}_{\lambda (x)}(x)\) ，由于 ad \(x\) 半单。

推论 2. g 的每个正则元半单。

事实上，这样的元素属于某个嘉当子代数（第 3 节，定理 1 (i)）。

推论 3. 设 \(\mathfrak{h}\) 是约化李代数 \(\mathfrak{g}\) 的嘉当子代数。a) \(\mathfrak{h}\) 交换。

\begin{enumerate}
\def\labelenumi{\alph{enumi})}
\setcounter{enumi}{1}
\tightlist
\item
  若 \(\rho\) 是 \(\mathfrak{g}\) 的一个有限维半单表示，则 \(\rho(\mathfrak{h})\) 的元素半单。
\end{enumerate}

设 \(\mathfrak{c}\) 是 \(\mathfrak{g}\) 的中心，\(\mathfrak{s}\) 是它的导代数。那么 \(\mathfrak{g} = \mathfrak{c} \times \mathfrak{s}\) ，故 \(\mathfrak{h} = \mathfrak{c} \times \mathfrak{h}'\) ，其中 \(\mathfrak{h}'\) 是 \(\mathfrak{s}\) 的嘉当子代数（命题 2）。注意到定理 2，\(\mathfrak{h}'\) 交换，故 \(\mathfrak{h}\) 亦然。此外，\(\rho(\mathfrak{h}')\) 由半单元素组成，\(\rho(\mathfrak{c})\) 亦然（第一章，§6，第 5 节，定理 4）；由此得到断言 \(b\) )。

